\section{Cluster expansion}\label{sec:cluster}
%
On the surface our proof of Theorem~\ref{thm:alg,cluster} has a similar flavour to the polynomial interpolation method: we define a series expansion for $\log Z(G;w)$, show that it converges, and approximate $Z(G;w)$ by computing the coefficients of a truncation of the series. 
Instead of working with a Taylor series and a zero-free region, we work with a different formal power series for $\log Z$ called the \emph{cluster expansion} which expresses $\log Z$ as a sum of some weights over connected subgraphs of $G$. 
This technique was recently applied to approximating the partition functions of the Potts and random cluster models in~\cite{BCHPT19}, where the \emph{random cluster model} is random graph model from statistical physics that generalizes the Ising and Potts models, and percolation\footnote{Note the distinct uses of the term `cluster' in `cluster expansion' and `random cluster model', though there is a common theme of connected subgraphs in both uses.}.
To obtain the result we adapt a standard reduction to express our partition function $Z(G;w)$ in terms of the random cluster model, and apply the method of~\cite{BCHPT19} which gives an approximation algorithm via the cluster expansion.
%
\subsection{The random cluster model}
%
The random cluster model, instead of counting graph labelings according to satisfied edges, counts connected subgraphs according to some weights. 
We adapt the standard reduction comparing the Potts model and random cluster model partition functions to our $Z(G;w)$ for Unique Games instances.
%
We start by rewriting $Z(G;w)$ for a given instance $G=(V,E,\pi)$. 
Let $(V',E')$ be a connected component of $G$, so that $V'$ is a nonempty subset of $V$ and $E'\subset E\cap \binom{V'}{2}$. 
We consider a singleton vertex $\{u\}$ to comprise the connected component $(\{u\},\emptyset)$. 
Define
\[
\sat_\pi(V',E'):=\sum_{\{x_u\}_{u\in V'}\in[k]^{V'}}\;\prod_{(u,v)\in E'}\indicator{x_v=\pi_{uv}(x_u)},
\]
where $\indicator{P}=1$ if $P$ is true, and $0$ otherwise. 
In other words $\sat_\pi(V',F')$ counts the number of assignments of value $1$ (perfectly satisfying assignments) of the Unique Games instance restricted to the subgraph $(V',E')$. 
The definition means that $\sat_\pi(\{u\},\emptyset)=k$ as there are no constraints and the empty product is $1$.
Since we work with $(V',E')$ connected, we also have 
\begin{equation}\label{eq:con}
0\le \sat_\pi(V', E')\le k\,,
\end{equation}
as given any starting color for an arbitrary first vertex $u\in V'$, there is at most one completion of the coloring to a perfectly satisfying assignment obtained by following the constraints out along the component from $u$. 
%
We use the notation $\comp(V, F)$ for the set of connected components of the graph $(V, F)$, taken as pairs $(V', E')$ with $E'\subset F$. 
The following lemma gives the reduction from $Z(G;w)$ to the random cluster model partition function.
%
%
\begin{restatable}{lemma}{randomcluster}\label{lem:random cluster}
Let $G=(V,E,\pi)$ be a UG instance and $w\in\CC$.
Then 
\[
Z(G;w)=\sum_{F\subseteq E} (w-1)^{|F|} \prod_{(V',E')\in\comp(V,F)} \sat_{\pi}(V',E').
\]
\end{restatable}
%

%
\subsection{The cluster expansion}
%
We closely follow the notation and setup of~\cite{BCHPT19,HPR19}. Given a UG instance $G=(V,E,\pi)$, define a \emph{polymer} $\gamma$ to be a connected subgraph of $G$ with at least two vertices. A collection of polymers is \emph{compatible} if the polymers contained in it are pairwise vertex disjoint.
We define the incompatibility graph $H_G$ on the collection of all polymers as follows: vertices of $H_G$ are the polymers and two polymers are connected by an edge if they are not compatible (that is if they share a vertex). 
Write $|\gamma|:=|V(\gamma)|$, $\|\gamma \|:=|E(\gamma)|$, and given $w\in\CC$, define the \emph{weight} of a polymer $\gamma$ as 
\[
w_\gamma:=(w-1)^{\|\gamma\|}k^{-|\gamma|}\sat_\pi(\gamma)\,,
\]
where we write $\sat_\pi(\gamma)$ for the more cumbersome $\sat_\pi(V(\gamma),E(\gamma))$.
Then by Lemma~\ref{lem:random cluster} and the observation that for a single vertex $u$ we have $\sat_\pi(\{u\},\emptyset)=k$, we have 
\[
\Xi(G):=\sum_{\Gamma=\{\gamma_1,\ldots,\gamma_t\}}\prod_{i=1}^t w_{\gamma_i}=k^{-|V|}Z(G;w),
\]
where the sum is over all sets $\Gamma$ of (pairwise) compatible polymers. 
Note that $\Xi(G)$ is the multivariate independence polynomial of the compatibility graph $H_G$.
%
%
The \emph{cluster expansion} is the following formal power series for $\log\Xi(G)$:
\begin{equation}\label{eq:cluster expansion}
\log\Xi(G)=\sum_{\substack{\Gamma\subset V(H_G)\\\mathclap{H_G[\Gamma] \text{ connected}}}}\phi(\Gamma)\;\prod_{\gamma\in \Gamma}w_\gamma,
\end{equation}
where $\phi(\Gamma)$ is the \emph{Ursell function} of the graph $H_G[\Gamma]=(\Gamma,F)$, defined as
\[
\phi(\Gamma):=
\frac{1}{|\Gamma|!}\sum_{\substack{A\subseteq F\\ (\Gamma, A)\text{ connected}}} (-1)^{|A|}.
\]
For $\Gamma\subset V(H_G)$, let $\|\Gamma\|$ be given by $\|\Gamma\|:=\sum_{\gamma\in \Gamma}\|\gamma\|$, and define the truncated cluster expansion as follows
\begin{equation}
T_m:=\sum_{\substack{\mathclap{\Gamma\subset V(H_G),\, \|\Gamma\|< m}\\H_G[\Gamma] \text{connected}}}\phi(\Gamma)\;\prod_{\gamma\in \Gamma}w_\gamma\,.
\end{equation}
%
With the definitions and a reduction to the right partition function in place, we can now state the conditons of~\cite{BCHPT19} that imply the cluster expansion converges and gives an approximation guarantee.
%
\begin{restatable}[{Borgs et al.~\cite[Lemma2.1]{BCHPT19}}]{lemma}{KPconvergence}\label{lem:KPconvergence}
Suppose that polymers are connected subgraphs containing at least one edge of a graph $G$ of maximum degree $\Delta$ on $n$ vertices, that $\|\gamma\|$ is the number of edges of the polymer $\gamma$, and that
\begin{equation}|w_\gamma|\leq e^{-(7+\log\Delta)\|\gamma\|}.\label{eq:ce 2}\end{equation}
Then the cluster expansion converges absolutely and for any $m\in\mathbb N$, $|T_m-\log\Xi(G)|\leq ne^{-3m}$.
\end{restatable}

%
\noindent
To prove Theorem~\ref{thm:alg,cluster} we simply check that these conditions hold, which we state as a lemma below. 
%
\begin{restatable}{lemma}{clusterconvergence}\label{lem:clusterconvergence}
Let $\Delta\in \mathbb{N}_{\geq 16}$, let $C=e^{-9-2\log\Delta}$, and let $\zeta= 8\sqrt{1/\Delta}$.
Then if $k\geq C^{-2\Delta/\zeta}$ and $1\leq w\leq e^{(2-\zeta)\log(k)/\Delta}$, Lemma~\ref{lem:KPconvergence} holds for \UG[k] instances $G$ of maximum degree $\Delta$.
\end{restatable}
%

%
We deduce the following runtime guarantees from our setup and the analyses of~\cite{HPR19,PR17}.
The truncated series $T_m$ can be computed in time $e^{O(\Delta m+\log n)}$ given an enumeration of all polymers on fewer than $m$ edges and their weights (see~\cite{HPR19}). 
We can enumerate the polymers in time $O(n^2m^7(e\Delta)^{2m})$ as they are connected subgraphs of a graph of maximum degree $\Delta$ (see~\cite{PR17}), and compute each weight in time $O(km)$ as all perfectly satisfying assignments on a connected graph are found by following each of the $k$ assignments of an initial vertex and propagating along constraints.
To get an approximation of the form $e^{-\alpha}\le Z(G;w^*)/\xi\le e^\alpha$ we take $m=\log(n/\alpha)/3$ which means the entire computation of $\xi$ can be done in time
\[
e^{O(\Delta m+\log n)} + O(km^8n^2(e\Delta)^{2m})= kn^{O(1)}(n/\alpha)^{O(\Delta)}\,.
\]
We see that the number of colours needed to make the lemma work is of the order $\Delta^{O(\Delta^{3/2})}$. It would be interesting to get a better dependence on $\Delta$.
%
%
